% Generated draft: requires mathematical and source review.
\documentclass[11pt]{article}
\usepackage[margin=25mm]{geometry}
\usepackage{fontspec}
\setmainfont{Latin Modern Roman}
\usepackage{amsmath,amssymb,mathtools,unicode-math}
\setmathfont{Latin Modern Math}
\usepackage{enumitem,needspace,xurl,hyperref}
\setlength{\parindent}{0pt}
\setlength{\parskip}{5pt}
\setlength{\emergencystretch}{3em}
\newcommand{\N}{\mathbb{N}}
\newcommand{\Z}{\mathbb{Z}}
\newcommand{\Q}{\mathbb{Q}}
\newcommand{\R}{\mathbb{R}}
\newcommand{\C}{\mathbb{C}}
\newcommand{\F}{\mathbb{F}}
\newcommand{\A}{\mathbb{A}}
\newcommand{\PP}{\mathbb{P}}
\newcommand{\E}{\mathbb{E}}
\newcommand{\eps}{\varepsilon}
\newcommand{\dd}{\,\mathrm{d}}
\begin{document}
\section*{UnsolvedMath ID: 7200085}
\textbf{Source problem code:} catalogue entry
\subsection{Definitions and mathematical statement}
Let \(P\subset\R^3\) be a full-dimensional compact convex polyhedron and let \(\pi:\R^3\to\R^2\) be orthogonal projection to a fixed plane. In the standard strict-clearance formulation, \(P\) has Rupert's property when
\[
\exists R,S\in SO(3)\ \exists a\in\R^2,\qquad
 a+\pi(RP)\subset\operatorname{int}(\pi(SP)).
\]
The interior is planar interior. This is the projection formulation of passing a congruent copy through a straight hole cut through the polyhedron in a suitable orientation, with positive clearance; rotations and translations are permitted but no scaling is allowed.

The proposed universal assertion is
\[
\forall P\text{ full-dimensional compact convex polyhedron},
\qquad P\text{ has Rupert's property}.
\]
A counterexample is a single such \(P\) for which the displayed projection inclusion fails for every \(R,S,a\); testing a finite selection of orientations alone would not establish that failure.

The supplied research summary records a 2025 counterexample, arXiv:2508.18475, and labels the universal conjecture solved negatively. It dates the conjecture to 2017 and says the cube and some other prisms remain positive examples; these comments do not assert that all prisms have the property. The counterexample claim and its universal orientation obstruction are supplied literature metadata, not verified by this draft. General characterization is mentioned as an ongoing topic, separate from the yes/no question here.
\subsection{Short English statement}
Must every convex polyhedron admit one rotated projection fitting strictly inside another without scaling? This is Rupert's property; the supplied registry reports a 2025 counterexample to the universal claim.
\subsection{Sources}
\begin{itemize}
\item Local UnsolvedMath registry, permanent ID \(7200085\), exact source code \texttt{catalogue entry}, ``Does every convex polyhedron have Rupert's property.''
\item Wikipedia, \emph{List of unsolved problems in mathematics}, pinned revision 1366636767, Geometry bullet 85; \url{https://en.wikipedia.org/w/index.php?title=List_of_unsolved_problems_in_mathematics&oldid=1366636767#Geometry}. MediaWiki-source extraction; recorded access 2026-07-29.
\item Supplied research summary cites arXiv:2508.18475 (2025) for a convex-polyhedron counterexample. Authors, title, and publication details are not supplied, and no counterexample proof has been independently checked here.
\end{itemize}
\end{document}
